% GENERATED FILE. Edit canonical lessons, then run npm run generate:books.
% canonical-source-hash: fnv1a64:abd7a5b92af4e74a
% canonical-lessons: IT-C282-cronaca, IT-C282-quotidiano, IT-C282-settimanale, IT-C282-titolo, IT-C282-articolo

\chapter{Il giornale del mattino}
\label{ch:it-il-giornale-del-mattino}
\hlchaptermodality{282}

\begin{chapteropening}
\textbf{By the end of this chapter:} \emph{I can talk about the papers: the news pages, the dailies, the weeklies, a headline and an article.}

Complete the last lesson of chapter 282: I can talk about the papers: the news pages, the dailies, the weeklies, a headline and an article.
\end{chapteropening}

\section[la cronaca]{la cronaca — the news pages, the local news; the chronicle}

\label{lesson:IT-C282-cronaca}



\begin{quote}
Before the new one: say the Italian for last night, then the Italian for the whole morning.
\end{quote}

\subsection*{You'll want to know: la cronaca}
\textbf{la cronaca} — \textquotedblleft{}the news pages, the local news; the chronicle\textquotedblright{}.

\begin{quote}
Leggo sempre la cronaca della mia città.
\end{quote}

\begin{quote}
\emph{I always read the local news for my town.}
\end{quote}

\emph{La cronaca nera} is crime news; \emph{la cronaca rosa} is celebrity gossip.

\begin{cousinweb}
\emph{cronaca} is inherited from Latin \emph{chronica} \textquotedblleft{}a chronicle\textquotedblright{} (EtymDB), from Greek \emph{khronikós} \textquotedblleft{}of time\textquotedblright{}, built on \emph{khrónos} \textquotedblleft{}time\textquotedblright{}.
\end{cousinweb}

\begin{grammarlens}[title={What you've built}]
One more word for telling, explaining and describing at B1.
\end{grammarlens}

\subsection*{Your turn}
\begin{itemize}
\raggedright
  \item \emph{Say it:} \emph{la cronaca}
  \item \emph{Say it:} \emph{la cronaca}, once more
  \item \emph{Say it:} a sentence of your own with \emph{la cronaca}
  \item \emph{Recall:} say the Italian for last night, then the Italian for the whole morning, then say \emph{la cronaca} again
\end{itemize}

\begin{tcolorbox}[breakable,colback=teal!4,colframe=teal!35!black,title={Before you move on}]
What does la cronaca mean? (\textquotedblleft{}The news pages, the local news; the chronicle\textquotedblright{}.) Say it once more.
\end{tcolorbox}
\section[il quotidiano]{il quotidiano — the daily newspaper}

\label{lesson:IT-C282-quotidiano}



\begin{quote}
Before the new one: say the Italian for the whole morning, then the Italian for the news pages.
\end{quote}

\subsection*{You'll want to know: il quotidiano}
\textbf{il quotidiano} — \textquotedblleft{}the daily newspaper\textquotedblright{}.

\begin{quote}
Compro un quotidiano ogni mattina.
\end{quote}

\begin{quote}
\emph{I buy a daily paper every morning.}
\end{quote}

As an adjective, \emph{quotidiano} means daily, everyday: \emph{la vita quotidiana}.

\begin{cousinweb}
\emph{quotidiano} is a learned borrowing from Latin \emph{quotīdiānus} \textquotedblleft{}daily\textquotedblright{} (EtymDB), built on \emph{quot} \textquotedblleft{}as many as\textquotedblright{} and \emph{diēs} \textquotedblleft{}day\textquotedblright{}.
\end{cousinweb}

\begin{grammarlens}[title={What you've built}]
One more word for telling, explaining and describing at B1.
\end{grammarlens}

\subsection*{Your turn}
\begin{itemize}
\raggedright
  \item \emph{Say it:} \emph{il quotidiano}
  \item \emph{Say it:} \emph{il quotidiano}, once more
  \item \emph{Say it:} a sentence of your own with \emph{il quotidiano}
  \item \emph{Recall:} say the Italian for the whole morning, then the Italian for the news pages, then say \emph{il quotidiano} again
\end{itemize}

\begin{tcolorbox}[breakable,colback=teal!4,colframe=teal!35!black,title={Before you move on}]
What does il quotidiano mean? (\textquotedblleft{}The daily newspaper\textquotedblright{}.) Say it once more.
\end{tcolorbox}
\section[il settimanale]{il settimanale — the weekly (magazine or paper)}

\label{lesson:IT-C282-settimanale}



\begin{quote}
Before the new one: say the Italian for the news pages, then the Italian for the daily newspaper.
\end{quote}

\subsection*{You'll want to know: il settimanale}
\textbf{il settimanale} — \textquotedblleft{}the weekly (magazine or paper)\textquotedblright{}.

\begin{quote}
Mia madre legge un settimanale di cucina.
\end{quote}

\begin{quote}
\emph{My mother reads a weekly cookery magazine.}
\end{quote}

Built on \emph{settimana}, ‘week’.

\begin{cousinweb}
\emph{settimanale} is built on \emph{settimana} \textquotedblleft{}week\textquotedblright{}, from Late Latin \emph{septimāna}, built on \emph{septem} \textquotedblleft{}seven\textquotedblright{} (Lewis \& Short).
\end{cousinweb}

\begin{grammarlens}[title={What you've built}]
One more word for telling, explaining and describing at B1.
\end{grammarlens}

\subsection*{Your turn}
\begin{itemize}
\raggedright
  \item \emph{Say it:} \emph{il settimanale}
  \item \emph{Say it:} \emph{il settimanale}, once more
  \item \emph{Say it:} a sentence of your own with \emph{il settimanale}
  \item \emph{Recall:} say the Italian for the news pages, then the Italian for the daily newspaper, then say \emph{il settimanale} again
\end{itemize}

\begin{tcolorbox}[breakable,colback=teal!4,colframe=teal!35!black,title={Before you move on}]
What does il settimanale mean? (\textquotedblleft{}The weekly (magazine or paper)\textquotedblright{}.) Say it once more.
\end{tcolorbox}
\section[il titolo]{il titolo — the headline; the title}

\label{lesson:IT-C282-titolo}



\begin{quote}
Before the new one: say the Italian for the daily newspaper, then the Italian for the weekly.
\end{quote}

\subsection*{You'll want to know: il titolo}
\textbf{il titolo} — \textquotedblleft{}the headline; the title\textquotedblright{}.

\begin{quote}
Il titolo in prima pagina è enorme.
\end{quote}

\begin{quote}
\emph{The headline on the front page is huge.}
\end{quote}

In a paper \emph{i titoli} are the headlines; in a book, \emph{il titolo} is the title.

\begin{cousinweb}
\emph{titolo} is inherited from Latin \emph{titulus} \textquotedblleft{}an inscription, a label, a title\textquotedblright{} (EtymDB).
\end{cousinweb}

\begin{grammarlens}[title={What you've built}]
One more word for telling, explaining and describing at B1.
\end{grammarlens}

\subsection*{Your turn}
\begin{itemize}
\raggedright
  \item \emph{Say it:} \emph{il titolo}
  \item \emph{Say it:} \emph{il titolo}, once more
  \item \emph{Say it:} a sentence of your own with \emph{il titolo}
  \item \emph{Recall:} say the Italian for the daily newspaper, then the Italian for the weekly, then say \emph{il titolo} again
\end{itemize}

\begin{tcolorbox}[breakable,colback=teal!4,colframe=teal!35!black,title={Before you move on}]
What does il titolo mean? (\textquotedblleft{}The headline; the title\textquotedblright{}.) Say it once more.
\end{tcolorbox}
\section[l'articolo]{l'articolo — the article (in a paper)}

\label{lesson:IT-C282-articolo}



\begin{quote}
Before the new one: say the Italian for the weekly, then the Italian for the headline.
\end{quote}

\subsection*{You'll want to know: l'articolo}
\textbf{l'articolo} — \textquotedblleft{}the article (in a paper)\textquotedblright{}.

\begin{quote}
Ho letto un articolo interessante sul treno.
\end{quote}

\begin{quote}
\emph{I read an interesting article on the train.}
\end{quote}

In grammar, \emph{l'articolo} is the article too: \emph{il}, \emph{la}, \emph{un}.

\begin{cousinweb}
\emph{articolo} is a learned borrowing from Latin \emph{articulus} \textquotedblleft{}a small joint, a part\textquotedblright{} (Lewis \& Short), a little \emph{artus} \textquotedblleft{}joint\textquotedblright{}.
\end{cousinweb}

\begin{grammarlens}[title={What you've built}]
One more word for telling, explaining and describing at B1.
\end{grammarlens}

\subsection*{Your turn}
\begin{itemize}
\raggedright
  \item \emph{Say it:} \emph{l'articolo}
  \item \emph{Say it:} \emph{l'articolo}, once more
  \item \emph{Say it:} a sentence of your own with \emph{l'articolo}
  \item \emph{Recall:} say the Italian for the weekly, then the Italian for the headline, then say \emph{l'articolo} again
\end{itemize}

\begin{tcolorbox}[breakable,colback=teal!4,colframe=teal!35!black,title={Before you move on}]
What does l'articolo mean? (\textquotedblleft{}The article (in a paper)\textquotedblright{}.) Say it once more.
\end{tcolorbox}
